% Generated draft: requires mathematical and source review.
\documentclass[11pt]{article}
\usepackage[margin=25mm]{geometry}
\usepackage{fontspec}
\setmainfont{Latin Modern Roman}
\usepackage{amsmath,amssymb,mathtools,unicode-math}
\setmathfont{Latin Modern Math}
\usepackage{enumitem,needspace,xurl,hyperref}
\setlength{\parindent}{0pt}
\setlength{\parskip}{5pt}
\setlength{\emergencystretch}{3em}
\newcommand{\N}{\mathbb{N}}
\newcommand{\Z}{\mathbb{Z}}
\newcommand{\Q}{\mathbb{Q}}
\newcommand{\R}{\mathbb{R}}
\newcommand{\C}{\mathbb{C}}
\newcommand{\F}{\mathbb{F}}
\newcommand{\A}{\mathbb{A}}
\newcommand{\PP}{\mathbb{P}}
\newcommand{\E}{\mathbb{E}}
\newcommand{\eps}{\varepsilon}
\newcommand{\dd}{\,\mathrm{d}}
\begin{document}
\section*{UnsolvedMath ID: 30001352}
\textbf{Source problem code:} catalogue entry
\subsection{Definitions and mathematical statement}
The supplied record does not define the section space containing \(\sigma\), the construction of \(Y_\sigma\), or the class of surfaces containing \(S\). Consequently, the following records the requested constructions with these parameters unresolved.

For a surface \(S\) over a field \(k\), write
\[
S^{[2]}=\operatorname{Hilb}^{2}(S)
\]
for the Hilbert scheme of zero-dimensional subschemes of length two. The record does not specify \(k\) or impose an explicit hypothesis on \(S\); the K3 surfaces mentioned in reference titles do not by themselves supply such a hypothesis.

For a vector bundle \(E\) on a variety \(X\), global generation means that the evaluation morphism
\[
H^0(X,E)\otimes\mathcal O_X\longrightarrow E
\]
is surjective. Given an algebraic group \(G\) acting transitively on \(X\), a homogeneous vector bundle means a vector bundle equipped with a compatible \(G\)-action linear on fibers. For \(s\in H^0(X,E)\), let \(Z(s)\) denote its scheme-theoretic zero locus. In this algebraic setting, a property holds for a general section if it holds for every section in some nonempty Zariski-open subset of \(H^0(X,E)\).

Under the usual compact complex convention, an irreducible symplectic manifold is a connected, simply connected compact K\"ahler manifold \(M\) such that
\[
H^0(M,\Omega_M^2)=\C\omega,
\]
where \(\omega\) is an everywhere nondegenerate holomorphic two-form. The record does not explicitly specify this convention or the underlying field.

The source asks for the following.
\begin{enumerate}
\item Find specific sections \(\sigma\) for which there is an isomorphism
\[
Y_\sigma\cong S^{[2]}.
\]
The intended section space, construction of \(Y_\sigma\), allowed surfaces \(S\), and category of isomorphism require source clarification. A birational equivalence alone does not satisfy the stated isomorphism requirement.
\item Find other globally generated homogeneous vector bundles whose general zero locus is an irreducible symplectic manifold. For a specified admissible pair \((X,E)\), the general-section requirement has the form
\[
\exists\,\varnothing\ne U\subseteq H^0(X,E)\text{ Zariski open}\quad
\forall s\in U,\quad Z(s)\text{ is an irreducible symplectic manifold}.
\]
The admissible ambient varieties, groups, bundle ranks, dimensions of zero loci, and examples relative to which the bundles must be ``other'' are not supplied. The question requests examples; it does not explicitly require a complete classification.
\end{enumerate}
The supplied dated assessment reports partial progress through Debarre--Voisin degenerations and special examples with isomorphic or birational Hilbert-square models, while describing an explicit \(\sigma\) and a general bundle classification as remaining open. These are supplied literature qualifications, not independently verified conclusions.
\subsection{Short English statement}
Find explicit sections whose associated varieties are isomorphic to Hilbert squares of the intended surfaces, and find further globally generated homogeneous bundles with irreducible symplectic general zero loci. The record omits the construction of the varieties, the section space, the surface hypotheses, and the permitted homogeneous spaces. Its supplied partial-solution assessment requires independent review.
\subsection{Sources}
\begin{itemize}
\item Supplied local UnsolvedMath registry record, permanent ID \(30001352\), title: Hilbert-Square Zero Loci of Homogeneous Bundles; exact source problem code: \texttt{catalogue entry}. Collection: Oberwolfach Reports Open Problems. No individual proposer is identified; the supplied proposed year is 2009.
\item Supplied report citation: \emph{Complex Algebraic Geometry (2010)}, Oberwolfach Reports, pp.\ 2483--2544, DOI \texttt{10.4171/owr/2009/44}. \url{https://doi.org/10.4171/owr/2009/44}. The record attributes the problem to workshop open-problem discussions documented in this report.
\item Supplied literature reference: \emph{Hilbert squares and Debarre-Voisin varieties}. \url{https://jep.centre-mersenne.org/item/JEP_2020__7__653_0/}.
\item Supplied literature reference: \emph{Hilbert squares of genus 16 K3 surfaces}. \url{https://arxiv.org/abs/2510.26488}.
\item Supplied literature reference: original Oberwolfach report. \url{https://doi.org/10.4171/owr/2009/44}.
\item Literature checked at: \texttt{2026-08-21}. This is supplied dated curation, not independently verified. It labels the problem partially solved and states that Debarre--Voisin degenerations and special examples give isomorphic or birational Hilbert-square models, but that the requested explicit \(\sigma\) and a general bundle classification remain open. The background explicitly requests independent verification of time-sensitive status and cited claims.
\item The registry labels the statement \texttt{corrected\_verified} and reports that its mathematical content was checked against the official report, with notation, context, and question form normalized without changing the source-level claim. This is a supplied verification assertion; no source consultation or independent verification is claimed here.
\end{itemize}
\end{document}
